% This LaTeX was auto-generated from MATLAB code.
% To make changes, update the MATLAB code and export to LaTeX again.

\documentclass{article}

\usepackage[utf8]{inputenc}
\usepackage[T1]{fontenc}
\usepackage{lmodern}
\usepackage{graphicx}
\usepackage{color}
\usepackage{hyperref}
\usepackage{amsmath}
\usepackage{amsfonts}
\usepackage{epstopdf}
\usepackage[table]{xcolor}
\usepackage{matlab}
\usepackage[paperheight=795pt,paperwidth=614pt,top=72pt,bottom=72pt,right=72pt,left=72pt,heightrounded]{geometry}

\sloppy
\epstopdfsetup{outdir=./}
\graphicspath{ {./lab1_media/} }

\matlabmultipletitles

\begin{document}

\matlabtitle{ Задание 1}


\vspace{1em}
\begin{matlabcode}
a_2 = 9;
a_1 = 23;
a_0 = 15;
b_2 = 14;
b_1 = 6;
b_0 = 16;
\end{matlabcode}

\begin{par}
\begin{flushleft}
Рассмотрим ДУ$\dddot{y} +a_2 \ddot{y} +a_1 \dot{y} +a_0 y=b_2 \ddot{u} +b_1 \dot{u} +b_0 u$
\end{flushleft}
\end{par}

\begin{par}
\begin{flushleft}
Перепишем в более удобном варианте для создания структурной схемы в Simulink
\end{flushleft}
\end{par}

\begin{par}
$$y=\frac{1}{p}\left(b_2 u-a_2 y+\frac{1}{p}\left(b_1 u-a_1 y+\frac{1}{p}\left(b_0 u-a_0 y\right)\right)\right)$$
\end{par}

\begin{matlabcode}
simtime = 10;
plot_1 = sim("lab1_1.slx", simtime);
plot(plot_1.y.Time, plot_1.y.Data, LineWidth=2, LineStyle="-");
hold on;
plot(plot_1.u.Time, plot_1.u.Data, LineWidth=2, LineStyle="--");
xlim([-0.1 10.00]);
ylim([0.00 1.25]);
grid on;
legend("$y(t)$", "$u(t)$", Location="best", FontSize=14, Interpreter="latex");
xlabel("$t, c$", Interpreter="latex");
ylabel("$Signal$", Interpreter="latex");
set(gcf, Position=[100, 100, 1200, 600]);
set(gca, FontSize=14); 
hold off;
exportgraphics(gcf, 'images/plot_1.png', Resolution=200);
\end{matlabcode}
\begin{center}
\includegraphics[width=\maxwidth{100.95333667837431em}]{figure_0.eps}
\end{center}


\matlabtitle{Задание 2}

\begin{par}
\begin{flushleft}
Найдём передаточную функцию из нашего ДУ
\end{flushleft}
\end{par}

\begin{par}
$$(p^3 +a_2 p^2 +a_1 p+a_0 )y=(b_2 p^2 +b_1 p+b_0 )u$$
\end{par}

\begin{par}
$$y=\frac{b_2 p^2 +b_1 p+b_0 }{p^3 +a_2 p^2 +a_1 p+a_0 }[u]\Rightarrow W(p)=\frac{b_2 p^2 +b_1 p+b_0 }{p^3 +a_2 p^2 +a_1 p+a_0 }$$
\end{par}

\begin{par}
\begin{flushleft}
Если подставить коэффициенты: $W(p)=\frac{14p^2 +6p+16}{p^3 +9p^2 +23p+15}$
\end{flushleft}
\end{par}

\begin{par}
\begin{flushleft}
Построим математические модели:
\end{flushleft}
\end{par}

\begin{par}
\begin{flushleft}
Общий вид матричной записи:
\end{flushleft}
\end{par}

\begin{par}
$$W(p)=(C(Ip-A)^{-1} B+D)$$
\end{par}

\begin{itemize}
\setlength{\itemsep}{-1ex}
   \item{\begin{flushleft} Каноническая управляемая форма \end{flushleft}}
\end{itemize}

\begin{par}
$$\begin{array}{l}
A_y =\left(\begin{array}{ccc}
0 & 1 & 0\\
0 & 0 & 1\\
-\frac{a_0 }{a_3 } & -\frac{a_1 }{a_3 } & -\frac{a_2 }{a_3 }
\end{array}\right)=\left(\begin{array}{ccc}
0 & 1 & 0\\
0 & 0 & 1\\
-15 & -23 & -9
\end{array}\right),~~B_y =\left(\begin{array}{c}
0\\
0\\
1
\end{array}\right)\\
C_y =\left(\begin{array}{ccc}
\frac{b_0 }{a_3 } & \frac{b_1 }{a_3 } & \frac{b_2 }{a_3 }
\end{array}\right)=\left(\begin{array}{ccc}
16 & 6 & 14
\end{array}\right),~~D_y =[d]=[0]
\end{array}$$
\end{par}

\begin{par}
\begin{flushleft}
$a_3$ в нашем случае равно 1, $d=0$.
\end{flushleft}
\end{par}

\begin{itemize}
\setlength{\itemsep}{-1ex}
   \item{\begin{flushleft} Каноническая наблюдаемая форма \end{flushleft}}
\end{itemize}

\begin{par}
$$\begin{array}{l}
A_н =\left(\begin{array}{ccc}
0 & 0 & -\frac{a_0 }{a_3 }\\
1 & 0 & -\frac{a_1 }{a_3 }\\
0 & 1 & -\frac{a_2 }{a_3 }
\end{array}\right)=\left(\begin{array}{ccc}
0 & 0 & -15\\
1 & 0 & -23\\
0 & 1 & -9
\end{array}\right),~~B_н =\left(\begin{array}{c}
\frac{b_0 }{a_3 }\\
\frac{b_1 }{a_3 }\\
\frac{b_2 }{a_3 }
\end{array}\right)=\left(\begin{array}{c}
16\\
6\\
14
\end{array}\right)\\
C_н =\left(\begin{array}{ccc}
0 & 0 & 1
\end{array}\right),~~D_H =[d]=[0]
\end{array}$$
\end{par}

\begin{itemize}
\setlength{\itemsep}{-1ex}
   \item{\begin{flushleft} Каноническая диагональная форма \end{flushleft}}
\end{itemize}

\begin{par}
$$A_д =\left(\begin{array}{ccc}
\lambda_1  & 0 & 0\\
0 & \lambda_2  & 0\\
0 & 0 & \lambda_3 
\end{array}\right),~~B_д =\left(\begin{array}{c}
\beta_0 \\
\beta_1 \\
\beta_2 
\end{array}\right),~~C_д =\left(\begin{array}{ccc}
1 & 1 & 1
\end{array}\right)$$
\end{par}

\begin{par}
$$\begin{array}{l}
\textrm{Найдём}\;\textrm{коэффициенты}\;{\lambda }:\\
p^3 +9p^2 +23p+15=0\\
\lambda_1 =-1\\
\frac{p^3 +9p^2 +23p+15}{p+1}=p^2 +8p+15\Rightarrow \lambda_2 =-3,~~\lambda_3 =-5\\
\textrm{Выразим}\;{\beta }:\\
\frac{14p^2 +6p+16}{(p+1)(p+3)(p+5)}=\frac{\beta_1 }{p+1}+\frac{\beta_2 }{p+3}+\frac{\beta_3 }{p+5}\\
\beta_1 (p^2 +8p+15)+\beta_2 (p^2 +6p+5)+\beta_3 (p^2 +4p+3)=14p^2 +6p+16\\
\left\lbrace \begin{array}{c}
\beta_1 +\beta_2 +\beta_3 =14\\
8\beta_1 +6\beta_2 +4\beta_3 =6\\
15\beta_1 +5\beta_2 +3\beta_3 =16
\end{array}\right.
\end{array}$$
\end{par}

\begin{matlabcode}
A = [1, 1, 1;
     8, 6, 4;
     15, 5, 3];
B = [14; 6; 16];
X = A \ B
\end{matlabcode}
\begin{matlaboutput}
X = 3x1    
    3.0000
  -31.0000
   42.0000

\end{matlaboutput}

\begin{par}
$$\beta_1 =3,~~\beta_2 =-31,~~\beta_3 =42$$
\end{par}

\begin{par}
\begin{flushleft}
Получаем каноническую диагональную форму:
\end{flushleft}
\end{par}

\begin{par}
$$A_д =\left(\begin{array}{ccc}
-1 & 0 & 0\\
0 & -3 & 0\\
0 & 0 & -5
\end{array}\right),~~B_д =\left(\begin{array}{c}
3\\
-31\\
42
\end{array}\right),~~C_д =\left(\begin{array}{ccc}
1 & 1 & 1
\end{array}\right)$$
\end{par}

\begin{matlabcode}
lam_1 = -1;
lam_2 = -3;
lam_3 = -5;
beta_1 = 3;
beta_2 = -31;
beta_3 = 42;
\end{matlabcode}

\begin{par}
\begin{flushleft}
Моделируем:
\end{flushleft}
\end{par}

\begin{matlabcode}
simtime = 10;
plot_Tf = sim("lab1_2_Tf.slx", simtime);
plot_contr = sim("lab_1_2_contr.slx", simtime);
plot_obs = sim("lab_1_2_obs.slx", simtime);
plot_diag = sim("lab_1_2_diag.slx", simtime);

plot(plot_Tf.u.Time, plot_Tf.u.Data, Linewidth=2, Linestyle="--", Color="magent");
hold on;
plot(plot_Tf.y.Time, plot_Tf.y.Data, LineWidth=2.5, LineStyle="-", Color="blue");
hold on;
plot(plot_contr.y.Time, plot_contr.y.Data, LineWidth=2.5, LineStyle="-.", Color="#FEC114");
hold on;
plot(plot_obs.y.Time, plot_obs.y.Data, LineWidth=2.5, LineStyle="--", Color="#72A71F");
hold on;
plot(plot_diag.y.Time, plot_diag.y.Data, LineWidth=2.5, LineStyle=":", Color="#FF4E00");

xlim([-0.1 10.00]);
ylim([0.00 1.25]);
grid on;
legend("$u(t)$", "Transfer Function", "Controllable Form", "Observable Form", ...
    "Diagonal Form", Location="best", FontSize=14, Interpreter="latex");
xlabel("$t, c$", Interpreter="latex");
ylabel("$Signal$", Interpreter="latex");
set(gcf, Position=[100, 100, 1200, 600]);
set(gca, FontSize=12); 
hold off;
exportgraphics(gcf, 'images/plot_2.png', Resolution=200);
\end{matlabcode}
\begin{center}
\includegraphics[width=\maxwidth{100.95333667837431em}]{figure_1.eps}
\end{center}


\matlabtitle{Задание 3}

\begin{par}
\begin{flushleft}
\includegraphics[width=\maxwidth{65.73005519317611em}]{image_0}
\end{flushleft}
\end{par}

\begin{par}
\begin{flushleft}
\includegraphics[width=\maxwidth{65.93075765178124em}]{image_1}
\end{flushleft}
\end{par}

\begin{par}
\begin{flushleft}
Рассмотрим систему $A(p)y(t)=B(p)u(t)$
\end{flushleft}
\end{par}

\begin{par}
$$A=\left(\begin{array}{cc}
p+15 & p+4\\
p+2 & p+6
\end{array}\right),~~B=\left(\begin{array}{cc}
6 & 6\\
4 & 3
\end{array}\right)$$
\end{par}

\begin{par}
\begin{flushleft}
Передаточная матрица:
\end{flushleft}
\end{par}

\begin{par}
$$W(p)=A^{-1} (p)B(p)$$
\end{par}

\begin{par}
\begin{flushleft}
Найдём $A^{-1}$:
\end{flushleft}
\end{par}

\begin{par}
$$\begin{array}{l}
\det A(p)=(p+15)(p+6)-(p+4)(p+2)=(15+6-4-2)p+15\cdot 6-4\cdot 2=15p+82\\
A^{-1} (p)=\frac{1}{15p+82}\left(\begin{array}{cc}
p+6 & -p-4\\
-p-2 & p+15
\end{array}\right)\\
W(P)=\frac{1}{15p+82}\left(\begin{array}{cc}
p+6 & -p-4\\
-p-2 & p+15
\end{array}\right)\left(\begin{array}{cc}
6 & 6\\
4 & 3
\end{array}\right)=\frac{1}{15p+82}\left(\begin{array}{cc}
6(p+6)-4(p+4) & 6(p+6)-3(p+4)\\
-6(p+2)+4(p+15) & -6(p+2)+3(p+15)
\end{array}\right)=\\
=\frac{1}{15p+82}\left(\begin{array}{cc}
2p+20 & 3p+24\\
48-2p & 33-3p
\end{array}\right)=\left(\begin{array}{cc}
\frac{2p+20}{15p+82} & \frac{3p+24}{15p+82}\\
\frac{48-2p}{15p+82} & \frac{33-3p}{15p+82}
\end{array}\right)
\end{array}$$
\end{par}

\begin{matlabcode}
simtime = 20;
plot_3 = sim("lab_1_3.slx", simtime);
plot(plot_3.y_1.Time, plot_3.y_1.Data, LineWidth=2, LineStyle="-");
hold on;
plot(plot_3.y_2.Time, plot_3.y_2.Data, LineWidth=2, LineStyle="-");
hold on;
plot(plot_3.u_1.Time, plot_3.u_1.Data, Linewidth=2, Linestyle="--");
hold on;
plot(plot_3.u_2.Time, plot_3.u_2.Data, Linewidth=2, Linestyle="--");
xlim([-0.1 20.00]);
ylim([-1.1 1.25]);
grid on;
legend("$y_1(t)$", "$y_2(t)$", "$u_1(t)$", "$u_2(t)$", Location="best", FontSize=14, Interpreter="latex");
xlabel("$t, c$", Interpreter="latex");
ylabel("$Signal$", Interpreter="latex");
set(gcf, Position=[100, 100, 1200, 600]);
set(gca, FontSize=14); 
hold off;
exportgraphics(gcf, 'images/plot_3.png', Resolution=200);
\end{matlabcode}
\begin{center}
\includegraphics[width=\maxwidth{100.95333667837431em}]{figure_2.eps}
\end{center}


\matlabtitle{Задание 4}

\begin{par}
\begin{flushleft}
Рассмотрим систему:
\end{flushleft}
\end{par}

\begin{par}
$$\begin{array}{l}
\left\lbrace \begin{array}{c}
\dot{x} =Ax+Bu\\
y=Cx
\end{array}\right.\\
A=\left(\begin{array}{cc}
0 & -7\\
1 & -3
\end{array}\right),~~B=\left(\begin{array}{cc}
9 & 3\\
4 & 5
\end{array}\right),~~C=\left(\begin{array}{cc}
4 & 6\\
5 & 1
\end{array}\right)
\end{array}$$
\end{par}

\begin{par}
\begin{flushleft}
Перепишем в более удобном виде для создания структурной схемы:
\end{flushleft}
\end{par}

\begin{par}
$$\begin{array}{l}
\left(\begin{array}{c}
x_1 \\
x_2 
\end{array}\right)=\frac{1}{p}\left(\begin{array}{cc}
0 & -7\\
1 & -3
\end{array}\right)\left(\begin{array}{c}
x_1 \\
x_2 
\end{array}\right)+\frac{1}{p}\left(\begin{array}{cc}
9 & 3\\
4 & 5
\end{array}\right)\left(\begin{array}{c}
u_1 \\
u_2 
\end{array}\right)\\
\left(\begin{array}{c}
x_1 \\
x_2 
\end{array}\right)=\left(\begin{array}{c}
-7\frac{1}{p}x_2 \\
\frac{1}{p}x_1 -3\frac{1}{p}x_2 
\end{array}\right)+\left(\begin{array}{c}
9\frac{1}{p}u_1 +3\frac{1}{p}u_2 \\
4\frac{1}{p}u_1 +5\frac{1}{p}u_2 
\end{array}\right)\Rightarrow \left\lbrace \begin{array}{c}
x_1 =\frac{1}{p}\left(-7x_2 +9u_1 +3u_2 \right)\\
x_2 =\frac{1}{p}\left(x_1 -3x_2 +4u_1 +5u_2 \right)
\end{array}\right.
\end{array}$$
\end{par}


\vspace{1em}
\begin{matlabcode}
simtime = 20;
plot_4 = sim("lab_1_4.slx", simtime);
plot(plot_4.y_1.Time, plot_4.y_1.Data, LineWidth=2, LineStyle="-");
hold on;
plot(plot_4.y_2.Time, plot_4.y_2.Data, LineWidth=2, LineStyle="-");
hold on;
plot(plot_4.u_1.Time, plot_4.u_1.Data, Linewidth=2, Linestyle="--");
hold on;
plot(plot_4.u_2.Time, plot_4.u_2.Data, Linewidth=2, Linestyle="--");
xlim([-0.1 20.00]);
ylim([-20.5 23]);
grid on;
legend("$y_1(t)$", "$y_2(t)$", "$u_1(t)$", "$u_2(t)$", Location="best", FontSize=14, Interpreter="latex");
xlabel("$t, c$", Interpreter="latex");
ylabel("$Signal$", Interpreter="latex");
set(gcf, Position=[100, 100, 1200, 600]);
set(gca, FontSize=10); 
hold off;
exportgraphics(gcf, 'images/plot_4.png', Resolution=200);
\end{matlabcode}
\begin{center}
\includegraphics[width=\maxwidth{100.95333667837431em}]{figure_3.eps}
\end{center}

\end{document}
